\documentclass{beamer}
\usepackage[utf8]{inputenc}
\usepackage[english,german]{babel} 
\usepackage{listings} 
\usepackage{listings-golang}

\usepackage[absolute,overlay]{textpos}
  \setlength{\TPHorizModule}{1mm}
  \setlength{\TPVertModule}{1mm}

\titlegraphic{\includegraphics[scale=0.1]{golanglogo.png}}
\title{
	\Large{\textit{\\Die Programmiersprache Go - Eine Einf\"uhrung}} \\
	\Large{\textbf{\\Konzeptskizze}}
}
\author{Student: Adrian Helberg \\Prüfer: Prof. Dr. Axel Schmolitzky}
\date{\today}

\definecolor{mygreen}{rgb}{0,0.6,0}

\begin{document}
\lstset{
    frame=single,
    basicstyle=\footnotesize,
    keywordstyle=\color{blue},
    showstringspaces=false, 
    stringstyle=\color{mygreen},
    tabsize=4,
    language=Golang
}

\maketitle

\frame{\tableofcontents}

\section{Geschichte}
\begin{frame}
\frametitle{Geschichte}

\begin{itemize}
\setlength{\itemsep}{30pt}
\item Entwickler
\item Entwurfsphase
\item Ver\"offentlichung
\item Go-Community
\end{itemize}

\end{frame}

\section{\"Uberblick}
\begin{frame}
\frametitle{\"Uberblick}

\begin{itemize}
\setlength{\itemsep}{20pt}
\item Motivation
\item Zielsetzung
\item Skalierbarkeit
\item Einsatzgebiete
\item Erster Vergleich mit Java\texttrademark
\end{itemize}

\begin{textblock}{20}(66,64)
\includegraphics[scale=0.1]{javalogo.png}
\end{textblock}

\end{frame}

\section{Merkmale}
\begin{frame}
\frametitle{Merkmale}

\begin{itemize}
\setlength{\itemsep}{16pt}
\item Closures
\item Reflexion
\item Typsicherheit
\item Autmoatische Speicherbereinigung
\item Interfaces, Mixins $\rightarrow$ Objektorientierung
\item Keine Klassen $\rightarrow$ Java Vergleich
\item Pakete
\end{itemize}

\end{frame}

\begin{frame}

\frametitle{Projekte}

\begin{textblock}{20}(10,30)
\textit{Go}
\includegraphics[scale=0.22]{docker.png}
\end{textblock}

\begin{textblock}{20}(60,32)
\textit{Java}
\includegraphics[scale=0.22]{elasticsearch.png}
\end{textblock}

\end{frame}

\section{Syntax}
\begin{frame}[fragile]
\frametitle{Syntax}

%%% GO
Go
\begin{lstlisting}
package main

import "fmt"

func main() {
    fmt.Println("Hello World!")
}
\end{lstlisting}

%%% JAVA
Java
\lstset{language=Java}
\begin{lstlisting}
public class HelloWorld 
{ 
    public static void main (String[] args)
    {
        System.out.println("Hello World!");
    }
}
\end{lstlisting}

\end{frame}

\section{Paradigmen}
\begin{frame}
\frametitle{Paradigmen}

\begin{itemize}
\setlength{\itemsep}{20pt}
\item Nebenläufig
\item Imperativ
\item Strukturiert
\item Modular
\item Objektorientiert
\end{itemize}

\end{frame}

\section{Zentrale Fragestellung}
\begin{frame}
\frametitle{Zentrale Fragestellung}

\centering
\glqq flexibel wie dynamisch getyped,\\ aber mit statischer Typsicherheit?\grqq{}

\end{frame}

\section{Pro \& Contra}
\begin{frame}
\frametitle{Pro \& Contra}

\begin{tabular}{p{5cm} | p{5.5cm}}
\textbf{Pro} & \textbf{Contra} \\ \hline
\begin{itemize}
\setlength{\itemsep}{20pt}
\item Minimalismus
\item Statisches Duck-Typing
\item Parallelisierung
\item Aufger\"aumte Syntax
\item Schneller Compiler
\end{itemize}
&
\begin{itemize}
\setlength{\itemsep}{20pt}
\item Keine generische Programmierung
\item nil statt Option
\item Wenig grundlegende Datenstrukturen
\item Keine Methodenüberladung
\end{itemize}
\\
\end{tabular}

\end{frame}

\begin{frame}
\frametitle{Pro \& Contra}

\begin{tabular}{p{5cm} | p{5.5cm}}
\textbf{Pro} & \textbf{Contra} \\ \hline
\begin{itemize}
\setlength{\itemsep}{20pt}
\item Statisch gelinkte Binärdateien
\item Laufzeiteigenschaften
\item Integriertes Unit-Test-Framework
\item Paketmanager
\end{itemize}
&
\begin{itemize}
\setlength{\itemsep}{20pt}
\item Unbefriedigende API-Dokumentation
\item Teilweise umständliche APIs
\item Umst\"andliches Mocking
\item Kleines \"Okosystem
\end{itemize}
\\
\end{tabular}

\end{frame}

\section{Compiler}
\begin{frame}
\frametitle{Compiler}

\begin{itemize}
\setlength{\itemsep}{20pt}
\item Gc
\item Gccgo
\end{itemize}

\end{frame}

\section{Einzelnachweise}
\begin{frame}
\frametitle{Einzelnachweise}

\begin{itemize}
\setlength{\itemsep}{20pt}
\item ...
\end{itemize}

\end{frame}

\end{document}